\subsubsection{LCS / Edit distance}

\paragraph{Idea intuitiva.}
\textbf{LCS} (Longest Common Subsequence): para dos strings $a$, $b$, $dp[i][j]$ = LCS de $a[0..i)$ y $b[0..j)$. Recurrencia:
\[ dp[i][j] = \begin{cases} dp[i-1][j-1] + 1 & \text{si } a[i-1] = b[j-1] \\ \max(dp[i-1][j], dp[i][j-1]) & \text{si no} \end{cases} \]

\textbf{Edit distance} (Levenshtein): numero minimo de operaciones (insert, delete, replace) para convertir $a$ en $b$. Recurrencia similar con un $+1$ adicional. La idea: ``alinear'' los strings eligiendo matches/mismatches/insert/delete.

\paragraph{Propiedades clave (invariantes).}
\begin{itemize}\setlength\itemsep{0pt}
	\item LCS y Edit Distance son DP $O(nm)$.
	\item Memoria se puede reducir a $O(\min(n, m))$ con rolling array.
	\item Edit distance con peso por operacion distinto: ajustar el \texttt{+1} por el peso.
	\item La subestructura optima: la respuesta para prefijos se construye de prefijos menores.
\end{itemize}

\paragraph{Codigo clave.}
La recurrencia clasica de LCS:
\begin{verbatim}
vector<vector<int>> dp(n + 1, vector<int>(m + 1, 0));
for (int i = 1; i <= n; ++i)
    for (int j = 1; j <= m; ++j)
        if (a[i-1] == b[j-1]) dp[i][j] = dp[i-1][j-1] + 1;
        else dp[i][j] = max(dp[i-1][j], dp[i][j-1]);
return dp[n][m];
\end{verbatim}

\paragraph{Complejidad.}
\begin{itemize}\setlength\itemsep{0pt}
	\item $O(NM)$ tiempo y memoria.
	\item Para $N, M \le 5000$: $\sim 2.5 \cdot 10^7$ operaciones.
\end{itemize}

\paragraph{Donde tocar para modificar.}
\begin{itemize}\setlength\itemsep{0pt}
	\item \textbf{Reconstruir la LCS / las operaciones}: aniadir \texttt{prev[i][j]} para backtrack.
	\item \textbf{LCS de varios strings}: NP-hard para $\ge 3$ strings; en la practica se hace con bitmask.
	\item \textbf{Edit distance con costos distintos}: cambiar el \texttt{+1} por el costo apropiado.
	\item \textbf{Rolling array}: si memoria es problema, mantener solo 2 filas.
	\item \textbf{Hirschberg}: algoritmo $O(NM)$ tiempo, $O(\min(N,M))$ memoria, con reconstruccion.
\end{itemize}

\paragraph{Trampas y errores frecuentes.}
\begin{itemize}\setlength\itemsep{0pt}
	\item Indices en la recurrencia: \texttt{a[i-1]} y \texttt{b[j-1]} (no \texttt{a[i]}).
	\item Para reconstruccion, el caso \texttt{a[i-1] == b[j-1]} tiene prioridad.
	\item Si los strings son muy largos, considerar Hirschberg o bitset optimization.
\end{itemize}

\paragraph{Codigo.}
Ver \texttt{cpp.tex}, seccion \textit{LCS / Edit distance}.

\vspace{0.5em}
